\documentclass{beamer}
\usefonttheme{professionalfonts}
%\usetheme{Rochester}
%\usetheme{AnnArbor}
\usetheme{Szeged}
\useoutertheme{infolines}
\usecolortheme{dolphin}
\usepackage{graphicx}
% \usepackage[english]{babel}
\usepackage{hyperref}
\usepackage{multicol}
\usepackage{color}
\usepackage{xcolor}
%\usepackage{smartdiagram}


%\definecolor{carmine}{rgb}{0.59, 0.0, 0.09}
\usepackage{amsmath} 

\usepackage{graphicx}
\usepackage{tikz}
\usetikzlibrary{shapes.callouts}
\usepackage{multirow}

\usepackage[utf8]{inputenc}

\usepackage{enumerate}
\usepackage[shortlabels]{enumitem}
\usepackage{amsmath,amsthm,amssymb,amsfonts,mathrsfs,latexsym,stmaryrd}
\usepackage{tcolorbox}

\usepackage{color} 
\usepackage{amssymb, amsmath, amsbsy}
\usepackage{adjustbox,lipsum}
\usepackage{placeins}

\usepackage{textpos}
\usepackage{wrapfig}
\usepackage[spanish, es-tabla]{babel}

\usepackage{caption}
\setlength{\parindent}{0pt}
\usepackage[hyperfootnotes=false]{hyperref}

\usepackage{url}
\usepackage{float}
\usepackage{tikz}
\usetikzlibrary{patterns}


\title[ICS-162]{Econometría ICS-294\\ \small{Clase 16: Heterocedasticidad }}

\author{Jocelyn Tapia Stefanoni}
\institute[USM]
{Universidad Técnica Federico Santa María\\ % Your institution for the title page
	\medskip
	\textit{} % Your email address
	
}
\logo{

\includegraphics[width=3cm]{figuras/logoUSM.png}
}
\date{}
\newcommand\Fontvi{\fontsize{9}{7.2}\selectfont}
\begin{document}
	
\begin{frame}
	\titlepage % Print the title page as the first slide


\end{frame}
%%%%%%%%%%%%%%%%%%%%%%%%%%%%%%%%%%%%%%%%%%%%%%%%%%%%%%%%%%%%%%%%%%%%%%%%%%%
\logo{}



\begin{frame}{Homocedasticidad vs Heterocedasticidad}
Los errores $u_i$ son heterocedásticos si la varianza condicional de $u_i$ dado $X_i$, $Var[u_i|X_i=x]$ no es constante.

\begin{figure}
    \centering
    \includegraphics[width=0.8\linewidth]{sem12/d2.png}
\end{figure}
   
\end{frame}

\begin{frame}{Consecuencias de la heterocedastidad}

\begin{itemize}
    \item El estimador MCO sigue siendo insesgado y consistente, $E[\hat{\beta}]$.
    \item Pero el estimador de MCO no es MELI.
    \item Los estimadores de las varianzas, $Var(\hat{\beta})$, serían sesgados.
    \item También, los estadísticos t no distribuirán t, ni los estadísticos F distribuirán F. 
    \item Estos problemas no se resuelven con tamaños de muestra grande.
\end{itemize}

\end{frame}



\begin{frame}{Ejemplo}
Considere el siguiente modelo lineal:

\begin{equation}
    Salario_i = \beta_0 + \beta_1 \cdot educacion_i + u_i
\end{equation}


En este contexto, analizamos si se cumple la \textbf{suposición de homocedasticidad} del término de error \( u_i \). Esta condición implica que la varianza del término de error, condicional a los años de educación, sea constante:

\begin{itemize}
    \item Para que haya homocedasticidad, debe cumplirse que: 
    \[
        \text{Var}[u_i \mid educacion_i] = \sigma^2
    \]
    \item Es decir, la varianza del término de error no debe depender del número de años de educación.
    \item Bajo homocedasticidad, los errores tienen varianza constante y no sistemáticamente relacionada con los niveles de educación.
\end{itemize}
 

    
\end{frame}

\begin{frame}{Ejemplo}
    \begin{figure}
        \centering
        \includegraphics[width=0.9\linewidth]{sem12/d5.png}
    \end{figure}
\end{frame}



\begin{frame}{Test para la heterocedasticidad (Breusch-Pagan)}

{\footnotesize
\begin{itemize}
    \item Considere el modelo MRL múltiple:
    $$y=\beta_0+\beta_1x_1+\beta_2x_2+...+\beta_kx_k+u$$
    \item La hipótesis nula de homocedasticidad es: 

    $$H_0:E(u^2|x_1,...,x_k)=E(u^2)=\sigma^2$$
    \item Se busca probar si $u^2$ se relaciona con las variables explicativas.
    \item Se puede suponer la siguiente relación:

    $$u^2=\delta_0+\delta_1x_1++...+\delta_kx_k+v$$
    \item Con la hipótesis nula:

    $$H_0: \delta_1=\delta_2=...=\delta_3=0$$

    \item No se conocen los errores de MLR original, pero se usan los residuos $\hat{u_i}$. Se estima:

    $$\hat{u}^2_i= \delta_0 +\delta_1x_1+...+\delta_kx_k+v$$

    \item Uitilizamos el test F de significancia global.
\end{itemize}
  }
\end{frame}

\begin{frame}{Test para la heterocedasticidad (Breusch-Pagan)}
    \begin{figure}
        \centering
        \includegraphics[width=0.9\linewidth]{sem12/d1.png}
    \end{figure}
\end{frame}

\begin{frame}{Test para la heterocedasticidad (Breusch-Pagan)}
    \begin{figure}
        \centering
        \includegraphics[width=0.9\linewidth]{d8.png}
    \end{figure}
\end{frame}

\begin{frame}{Test de White para la heterocedasticidad}
    \begin{figure}
        \centering
        \includegraphics[width=0.9\linewidth]{sem12/d9.png}
    \end{figure}

    La
prueba de White para heterocedasticidad es el estadístico ML y también puede utilizarse una prueba F.
\end{frame}

\begin{frame}{Test de White simplificado}
\begin{figure}
        \centering
        \includegraphics[width=0.9\linewidth]{sem12/d10.png}
    \end{figure}   

    Se
utilizan los valores ajustados porque éstos son funciones de las variables independientes.
\end{frame}



% \begin{frame}{Test ML}

%     El nombre estadístico del multiplicador de Lagrange proviene de la optimización con restricciones.

%     La forma del estadístico ML que se obtiene aquí se apoya en los supuestos de Gauss-Markov,
% los mismos supuestos que justifican el estadístico F en muestras grandes. El supuesto de normalidad
% no es necesario.

% Para obtener el estadístico ML considere el modelo de regresión múltiple usual con k variables
% independientes:
% $$y=\beta_0+\beta_1x_1+...+\beta_kx_k+u$$
% Se desea probar, por ejemplo, si las q últimas variables tienen parámetros poblacionales igual a
% cero, esto impone q restricciones de exclusión al modelo.

% \end{frame}

% \begin{frame}{Test ML}
%  El estadístico ML sólo requiere la estimación del modelo restringido.


%  Si las q variables omitidas  en realidad tienen coeficientes poblacionales iguales a
%  cero, entonces, por lo menos de manera aproximada, los residuos de la regresión restringida no deberían estar correlacionados con ninguna
%  de estas variables de la muestra. Para testearlo se corre una regresión auxiliar de dichos residuos sobre los k regresores. 

%  \begin{eqnarray*}    y=\tilde{\beta}_0+\tilde{\beta}_1x_1+...+\tilde{\beta}_{k-q}x_{k-q}+\tilde{u}\\
% \tilde{u}=\hat{\delta}_0+\hat{\delta}_1x_1+...+\hat{\delta}_kx_k
%  \end{eqnarray*}

%  La R-cuadrada de  debe tener un valor cercano a cero, asi $\tilde{u}$  no está correlacionada aproximadamente con ninguna de las variables independientes.
% \end{frame}

% \begin{frame}{Test ML}
    

% Como siempre en las pruebas de hipótesis, la cuestión es cómo determinar cuándo es
% suficientemente grande el estadístico para rechazar la hipótesis nula al nivel de significancia elegido.
% Resulta que, bajo la hipótesis nula, el tamaño de la muestra multiplicado por la R-cuadrada
% usual de la regresión auxiliar  tiene una distribución asintótica como variable aleatoria $\chi^2_q$.

%  \end{frame}
\begin{frame}{Inferencia robusta a la heterocedasticidad}
\begin{figure}
       \centering
       \includegraphics[width=0.9\linewidth]{sem12/d13.png}
   \end{figure}   
\end{frame}



% \begin{frame}{Recomendaciones}

% La forma funcional logarítmica de la variable dependiente suele reducir la heterocedasticidad.

% Si los residuales cuadrados se regresan
% sólo sobre una variable independiente, la prueba
% para la heterocedasticidad es precisamente
% el estadístico t usual de esa variable. Un estadístico
% t significativo sugiere que la heterocedasticidad
% es un problema.
    
% \end{frame}

% \begin{frame}{}
    
% \end{frame}
\end{document}